\documentclass{article}
\usepackage{amsmath}
\title{Continuous Stationary Dwell and Nonmonotone Read Feasibility}
\begin{document}\maketitle
\paragraph{Pin.} Use the same motion, two inspection requirements,
time horizon and initial rest as the acceleration-limited study.
A truthful notice becomes available immediately at the end of
continuous rest of length \(d\ge0\) at any safe position. It
completes at a prescribed time \(r\), with dwell beginning at
\(r-d\ge0\). There is no fixed station, dwell speed tolerance,
after-read latency or hardware guarantee. The same common
pre-read motion must work for both worlds.

\paragraph{Full feasible set.}
Write \(\sigma,k,\alpha\) for the far seam time, required minimum
seam speed, and acceleration magnitude. Let
\[
 s=k/\alpha,\qquad r_0=\sigma-s,\qquad
 D_*=\sigma-(1+\sqrt2)s.
\]
Then the exact condition is \(0\le d\le r\le r_0\) and
\[
 r-d\ \ge\
 \sqrt{\max\{0,\,2s^2-(r_0-r)^2\}}.
\]
To prove necessity, the far robot starts its suffix from rest,
reaches speed at least \(k\) in \(T=\sigma-r\), and at that time
has position at most 4. In particular \(T\ge s\).
With \(u=T-s\ge0\), its velocity at suffix time \(t\) is at least
\(\max(-\alpha t, k-\alpha(T-t))\). Integrating gives minimum
signed advance \(m=k^2/(2\alpha)-\alpha u^2/4\).
For \(m>0\) the read site must be no greater than \(4-m\).
Starting and ending the preparatory motion at rest bounds
its backward displacement in time \(r-d\) by
\(\alpha(r-d)^2/4\): its speed in magnitude is at most
the minimum of the two acceleration cones. This yields
the displayed root. For \(m\le0\) no retreat is forced.

For sufficiency when \(m>0\), move from 4 to \(4-m\) in
\(2\sqrt{m/\alpha}\) by a backward triangular velocity profile,
then hold through the dwell. When the notice is far, follow
the lower velocity envelope until \(\sigma\), arriving at
\((4,k)\), then follow the previously checked far suffix.
The pre-seam suffix first goes backward, then forward, never
crosses 4 before the seam, and its greatest retreat beyond
the read site is \(\alpha u^2/8\). During both the preparation
and suffix the second robot stays above \(4-\alpha s^2>1\).
Its pre-seam speed is at most \(k<1\), including the
first backward part as \(u\le\sqrt2s\). Both sensing footprints
therefore touch or overlap until the far seam; coverage on the
near portion follows from their far edge being above 3.
The late branch can instead hold indefinitely at the read
site, keeping the seam covered. If \(m\le0\), read at the
initial stationary position and hold until \(r=D_*\);
then use the far suffix for \(u=\sqrt2s\), where \(m=0\).
The late branch remains stationary. This deals with all earlier
cuts without assuming that an already chosen acquisition site
can be changed after the notice.

\paragraph{Consequences and replay.}
The maximum duration is \(D_*\): on the later segment
\(u+\sqrt{2s^2-u^2}\ge\sqrt2s\) for \(0\le u\le\sqrt2s\).
Both \(r=D_*\) and \(r=r_0\) attain \(D_*\), while the allowable
dwell at the middle cut \(r=\sigma-2s\) is only \(\sigma-3s\).
A fixed dwell of 0.42 is feasible at 0.425 and 0.548, but
not at 0.470. The latest completion for every feasible dwell
remains \(r_0\simeq0.548224\). The exact replay compares over
38,493 rational cases with independent trapezoid integration,
and certifies the stated examples using both endpoints of
a rational enclosure of \(k\). The duration threshold is
approximately \(0.435484\), not a measured sensing time.
An equally informed direct reference has the same feasible
set. The analysis does not transport to a fixed read site,
bounded nonzero speed, or a delayed notice.
\end{document}
